%---------------------------------------------------------------------
% Quick-reference sheet and the recurring mistakes.
%---------------------------------------------------------------------
\section*{فوری دہرائی}
\markboth{فوری دہرائی}{}

\begin{nukta}[چار فرق جو ہر بار کام آتے ہیں]
\textbf{شرط / رکن} --- نماز سے پہلے، باہر / نماز کے اندر، جزو۔ \quad
\textbf{فرض / واجب} --- چھوٹے تو نماز دوبارہ / چھوٹے تو سجدہ سہو۔

\smallskip
\textbf{غیبت / بہتان} --- عیب موجود ہو / موجود نہ ہو۔ \quad
\textbf{اتباع / اطاعت} --- قدم بہ قدم پیروی / حکم کی تعمیل۔
\end{nukta}

\medskip
\subsection*{سنین --- تاریخ کے سوال انہی پر جیتے جاتے ہیں}

\begin{center}
\begin{tabular}{@{}p{6.4cm} p{3.4cm} p{4.4cm}@{}}
\toprule
\textbf{واقعہ} & \textbf{ہجری} & \textbf{عیسوی} \\
\midrule
غزوہ احد                        & شوال \textenglish{3}ھ   & \textenglish{625}ء \\
خلافتِ امویہ                     & \textenglish{41--132}ھ  & \textenglish{661--750}ء \\
خلافتِ عباسیہ                    & \textenglish{132--656}ھ & \textenglish{750--1258}ء \\
ہارون الرشید                     & \textenglish{170--193}ھ & \textenglish{786--809}ء \\
فتحِ اندلس (طارق بن زیاد)        & \textenglish{92}ھ       & \textenglish{711}ء \\
\textbf{سقوطِ بغداد (ہلاکو)}     & \textbf{\textenglish{656}ھ} & \textbf{\textenglish{1258}ء} \\
سقوطِ غرناطہ                     & \textenglish{897}ھ      & \textenglish{1492}ء \\
\bottomrule
\end{tabular}
\end{center}

%---------------------------------------------------------------------
\subsection*{تین عام غلطیاں}

\begin{enumerate}\setlength{\itemsep}{3pt}
  \item \textbf{عربی غلط لکھنا۔} جتنی یاد ہے اتنی \emph{صحیح} لکھیے --- آدھی
        صحیح آیت پوری غلط آیت سے بہتر ہے۔
  \item \textbf{فرض، واجب اور شرط، رکن کو ایک سمجھنا۔} سجدہ سہو کا پورا مسئلہ
        اسی فرق پر کھڑا ہے --- خزاں \textenglish{2021} میں یہی سوال تھا۔
  \item \textbf{اموی و عباسی کارنامے ملا دینا؛ نمبر چالیس سمجھنا۔} سکہ و قبۃ
        الصخرہ \textbf{اموی}، بیت الحکمہ \textbf{عباسی}؛ اور ایسوسی ایٹ ڈگری
        میں کامیابی کا نمبر اب \textbf{پچاس} ہے۔
\end{enumerate}

% The closing note is two lines that keep landing alone on a nearly-blank
% final page: the preceding material fills page 21 to a hair short of
% where these two lines would fit. \enlargethispage borrows that hair's
% width back rather than leaving an entire page for two lines.
\enlargethispage{2\baselineskip}
\smallskip
\noindent
سب سے زیادہ یعنی \textbf{چاروں} پرچوں میں آئے موضوعات کی فہرست، اور ماخذ کی
تفصیل، کتابچے کے آغاز میں درج ہے۔
